% SPDX-FileCopyrightText: © 2026 Idiap Research Institute <contact@idiap.ch>
%
% SPDX-FileContributor: Andre Freitas <andre.freitas@idiap.ch>
% Neuro-symbolic AI Group
%
% SPDX-FileContributor: Danilo Gusicuma <danilo.gusicuma@idiap.ch>
% Neuro-symbolic AI Group
%
% SPDX-License-Identifier: MIT
\documentclass[11pt]{article}

\usepackage[margin=1in]{geometry}

\usepackage{amsmath,amssymb}

\title{Yang--Mills Big-Deal Goal Process Summary}

\author{Abducer Long-Horizon Monitor}

\date{}

\begin{document}

\maketitle

\section*{Archive Status}

Status: armed\_no\_reached\_goal. No milestone has been marked as reached yet.

\section*{Trigger Policy}

\begin{itemize}
\item Reached statuses: achieved, certified\_breakthrough, potential\_breakthrough\_signal, potentially\_achieved
\item When any success milestone enters a reached status, retain a LaTeX process summary and a paper-argument draft.
\item The paper draft must distinguish certified progress from a full Yang-Mills proof unless all Clay commitments are discharged in the same theory.
\end{itemize}

\section*{Watched Big-Deal Milestones}

\subsection*{Continuum Yang-Mills object constructed}

\begin{itemize}
\item Status: route\_in\_play\_not\_achieved
\item Archive triggered: False
\item Why this is a big deal: This would attack the existence half of the Clay target rather than only optimizing a regulated lattice model.
\item Potentially achieved when: A Balaban/RG-controlled cutoff sequence is shown to have a continuum limit with OS-positive Schwinger functions for compact simple gauge group Yang-Mills on R\textasciicircum{}4.
\item Success signal: Cutoff-removal theorem with a named limiting object S\_n and preserved gauge-invariant observables.
\item Next action: Formulate the exact continuum-passage lemma connecting Balaban-style estimates to OS-positive limiting Schwinger functions.
\end{itemize}

\paragraph{Certification Evidence}

\begin{itemize}
\item Uniform multiscale/RG bounds across the cutoff-removal sequence.
\item Tightness or convergence argument for the regulated measures or Schwinger functions.
\item Reflection positivity survives the limiting procedure.
\end{itemize}

\paragraph{Disqualifiers}

\begin{itemize}
\item Only finite-volume or fixed-lattice-spacing control is available.
\item The limit is not shown to be a quantum Yang-Mills object on R\textasciicircum{}4.
\item Gauge-group structure is lost or replaced by a scalar/free surrogate.
\end{itemize}

\subsection*{OS reconstruction yields the physical Hilbert space}

\begin{itemize}
\item Status: route\_in\_play\_not\_achieved
\item Archive triggered: False
\item Why this is a big deal: This would turn Euclidean data into the Hamiltonian arena where a mass gap has a precise spectral meaning.
\item Potentially achieved when: The limiting Schwinger functions satisfy the OS axioms strongly enough to reconstruct H\_YM, Omega, and H from the same continuum theory.
\item Success signal: A reconstruction packet names the Hilbert space, vacuum, Hamiltonian/semigroup, observable domain, and positivity argument.
\item Next action: Write the OS reconstruction obligation as a bridge lemma with inputs, theorem used, and output objects H\_YM, Omega, H.
\end{itemize}

\paragraph{Certification Evidence}

\begin{itemize}
\item Reflection positivity for the limiting Schwinger functions.
\item Euclidean covariance, symmetry, regularity, and clustering assumptions required by the reconstruction theorem.
\item Explicit identification of gauge-invariant observable classes carried into the Hilbert space.
\end{itemize}

\paragraph{Disqualifiers}

\begin{itemize}
\item The Hilbert space is postulated rather than reconstructed.
\item The reconstruction applies to a different regulated theory.
\item The observable class is too small to speak about the physical excitation spectrum.
\end{itemize}

\subsection*{Nontrivial interacting continuum theory witnessed}

\begin{itemize}
\item Status: route\_in\_play\_not\_achieved
\item Archive triggered: False
\item Why this is a big deal: This would prevent an apparent construction from collapsing into a Gaussian, empty, or purely classical limit.
\item Potentially achieved when: A gauge-invariant observable has limiting correlations or cumulants incompatible with the free/Gaussian theory.
\item Success signal: A nonzero connected correlation, Wilson-loop area/perimeter behavior, or higher cumulant survives cutoff removal.
\item Next action: Select one candidate gauge-invariant observable family and state the exact non-Gaussianity/nontriviality lemma it would need to satisfy.
\end{itemize}

\paragraph{Certification Evidence}

\begin{itemize}
\item Observable is defined in the same limiting continuum object.
\item Connected or higher-order correlation does not vanish in the limit.
\item The witness distinguishes the result from a free Gaussian fixed point.
\end{itemize}

\paragraph{Disqualifiers}

\begin{itemize}
\item The witness is only heuristic or simulation-level evidence.
\item The observable is not gauge invariant.
\item The nontriviality claim comes from a different approximation scheme than the constructed limit.
\end{itemize}

\subsection*{Cutoff-uniform Euclidean mass-gap witness obtained}

\begin{itemize}
\item Status: route\_in\_play\_not\_achieved
\item Archive triggered: False
\item Why this is a big deal: This would provide the first concrete spectral-gap witness rather than only a desired Hamiltonian inequality.
\item Potentially achieved when: Gauge-invariant Euclidean correlations decay exponentially with constants stable under continuum-limit passage.
\item Success signal: A family of observables O has |<O(t)O(0)>\_c| <= C\_O exp(-m t) with m > 0 independent of the removed cutoff.
\item Next action: Turn the correlator route into a decay-to-gap lemma with explicit uniform constants and observable-density assumptions.
\end{itemize}

\paragraph{Certification Evidence}

\begin{itemize}
\item The decay estimate is proved before and after taking the continuum limit.
\item The mass parameter m has a positive lower bound independent of cutoff scale.
\item The observable family is tied to the reconstructed physical Hilbert space.
\end{itemize}

\paragraph{Disqualifiers}

\begin{itemize}
\item Decay is available only at strong coupling or fixed lattice spacing.
\item Constants degenerate to zero under cutoff removal.
\item The observable family is not spectrally determining.
\end{itemize}

\subsection*{Positive Hamiltonian spectral gap proved}

\begin{itemize}
\item Status: route\_in\_play\_not\_achieved
\item Archive triggered: False
\item Why this is a big deal: This is the mass-gap half of the Clay target in the reconstructed physical Hilbert space.
\item Potentially achieved when: The reconstructed Hamiltonian has no non-vacuum spectrum in (0, Delta) for some Delta > 0.
\item Success signal: A theorem proves Spec(H restricted to P\_perp H\_YM) subset [Delta, infinity) for the same continuum theory.
\item Next action: Use the Hamiltonian derivation as a checkpoint: list every premise needed before E\_H((0, Delta)) P\_perp = 0 can be claimed.
\end{itemize}

\paragraph{Certification Evidence}

\begin{itemize}
\item H, Omega, and P\_perp come from OS/Wightman reconstruction of the constructed Yang-Mills theory.
\item A decay-to-gap or direct spectral argument excludes all non-vacuum states below Delta.
\item Gauge-invariant observables determine the relevant excitation sector.
\end{itemize}

\paragraph{Disqualifiers}

\begin{itemize}
\item The gap is only a finite-lattice transfer-matrix gap.
\item The spectral statement is asserted before Hilbert-space reconstruction.
\item The proof omits vacuum uniqueness or the projection onto non-vacuum excitations.
\end{itemize}

\subsection*{Same-theory Clay target integration achieved}

\begin{itemize}
\item Status: not\_achieved
\item Archive triggered: False
\item Why this is a big deal: This is the point where local successes stop being isolated and become a candidate solution architecture for the full problem.
\item Potentially achieved when: Existence, nontriviality, OS/Wightman adequacy, continuum limit, and positive mass gap are proved for one and the same Yang-Mills theory.
\item Success signal: A single dependency-closed proof packet connects all five essential commitments with no theory-switching or cutoff leakage.
\item Next action: Keep this as the top-level success gate and only promote lower milestones when they preserve all essential commitments.
\end{itemize}

\paragraph{Certification Evidence}

\begin{itemize}
\item All upstream milestones cite the same limiting object and observable algebra.
\item No open proof obligation blocks a required Clay commitment.
\item Failure tests for cutoff leakage, triviality, and insufficient reconstruction all pass.
\end{itemize}

\paragraph{Disqualifiers}

\begin{itemize}
\item Construction, nontriviality, and gap are proved in different theories.
\item Any essential target commitment is weakened without an accepted weakening rule.
\item The proof depends on an unvalidated heuristic or numerical premise.
\end{itemize}

\section*{Physics Mechanism Learning From Strategy Attempts}

\subsection*{First we tried to express the goal as a bounded Hamiltonian spectral-gap subclaim.}

\paragraph{Physics mechanism}

This is the clean energy-spectrum formulation of the Clay target. Once Euclidean Yang-Mills has produced OS-positive Schwinger functions, reconstruction should give the physical Hilbert space of gauge-invariant states, the vacuum vector, time translations, and a positive Hamiltonian. A mass gap means that after removing the vacuum line there are no arbitrarily soft physical excitations: the first glueball-like gauge-invariant energy lies at least Delta above the vacuum.

\paragraph{What did not pay off or remains blocked}

\begin{itemize}
\item This did not pay off as a proof route because it described the endpoint more than the construction mechanism.
\item X: the Hamiltonian formulation presupposes the reconstructed continuum Hilbert space rather than constructing it.
\item Y: the key inequality E\_H((0,Delta)) P\_perp = 0 is the hard mass-gap lemma itself, so naming it does not discharge it.
\item Z: a fixed-lattice transfer-matrix gap may disappear as the cutoff is removed, and gauge-invariant observables may fail to determine the full excitation spectrum.
\end{itemize}

\paragraph{What was learned physically}

\begin{itemize}
\item We learned that Hamiltonian language is indispensable as an endpoint test but weak as a discovery engine unless the reconstruction inputs are already available.
\item A claimed gap must always be asked backwards: where did H\_YM, Omega, H, P\_perp, and the observable algebra come from?
\item This converted the analysis from naming a desired spectral inequality into tracking the Euclidean construction and reconstruction prerequisites that make the inequality meaningful.
\end{itemize}

\paragraph{Decision}

Keep this as the mathematical endpoint and derivation checkpoint, but do not promote it as the primary proof route.

\subsection*{Then we moved to Euclidean correlator decay as a spectral-gap witness.}

\paragraph{Physics mechanism}

This is the Euclidean physics face of a gap. If a gauge-invariant observable O creates physical excitations from the vacuum, its connected Euclidean time correlation is a Laplace transform of the Hamiltonian spectral measure seen by O. Exponential decay with a cutoff-independent rate says that this spectral measure has no weight near zero. The route is physically meaningful only if the observables survive the continuum limit and are rich enough to see the vacuum-orthogonal excitation sector.

\paragraph{What did not pay off or remains blocked}

\begin{itemize}
\item It did not pay off as a standalone proof route because Euclidean decay can be true in a regulated or strong-coupling regime without proving a continuum Yang-Mills mass gap.
\item It left two hidden assumptions exposed: the decay constants must not vanish under cutoff removal, and the observable family must determine the whole physical excitation spectrum.
\item It also did not by itself supply nontriviality; an exponentially decaying sector is not enough if the limiting theory is Gaussian, empty, or not the constructed Yang-Mills object.
\end{itemize}

\paragraph{What was learned physically}

\begin{itemize}
\item Correlator decay is valuable only when it is a same-theory spectral witness: the observables, OS reconstruction, and Hamiltonian must all come from the same limiting Schwinger functions.
\item The right mathematical object to demand next is not just a decay estimate, but a cutoff-uniform decay-to-spectrum lemma with an observable-density hypothesis.
\item This taught the loop to treat strong-coupling lattice decay as evidence-task material until it carries a continuum-transfer certificate.
\end{itemize}

\paragraph{Decision}

Keep this route in the active proof queue, but require uniformity and spectral-determining observables before treating it as progress.

\subsection*{In parallel we moved to Balaban RG as the continuum-construction spine.}

\paragraph{Physics mechanism}

This is the scale-dynamics side. Starting with regulated gauge fields, the RG analysis must show that integrating out short-distance modes does not lose control of the non-Abelian interaction, does not break the gauge-invariant observable algebra, and has a continuum limit as the lattice spacing goes to zero and the volume goes to infinity. For the Clay problem, ultraviolet stability is not yet the explanation of the gap; it is the condition under which there is a nontrivial continuum vacuum whose correlations could carry a gap.

\paragraph{What did not pay off or remains blocked}

\begin{itemize}
\item It did not pay off as a mass-gap route by itself because ultraviolet/RG stability is not the same thing as an infrared spectral threshold.
\item It still leaves the OS/Wightman bridge open: a controlled effective action sequence must produce Schwinger functions with the axioms needed for reconstruction.
\item It can falsely look like progress if the loop celebrates multiscale control while never extracting a gauge-invariant decay or spectral witness.
\end{itemize}

\paragraph{What was learned physically}

\begin{itemize}
\item Balaban-style control should be used as the construction spine, not as a substitute for the mass-gap mechanism.
\item The decisive test is whether RG-controlled cutoff objects produce OS-positive continuum Schwinger functions and a non-Gaussian same-theory witness.
\item This shifted the loop toward an RG-to-OS bridge lemma instead of adding more renormalisation vocabulary.
\end{itemize}

\paragraph{Decision}

Keep this route in the active proof queue as the construction spine.

\subsection*{Now the recursive strategy is to fuse construction and spectral transfer.}

\paragraph{Physics mechanism}

The bridge has to say how a continuum gauge theory constructed by scale control becomes a spectral statement. The same limiting Schwinger functions must satisfy reflection positivity, define the gauge-invariant observable algebra, retain a non-Gaussian connected sector, and have Euclidean time correlations whose spectral measures avoid a neighborhood of zero. Only then can the reconstructed Hamiltonian inherit a positive lower threshold.

\paragraph{What did not pay off or remains blocked}

\begin{itemize}
\item It keeps the continuum object, OS reconstruction, nontriviality witness, and Hamiltonian gap in the same theory.
\item It prevents drift into pure Hamiltonian notation without construction.
\item It prevents drift into lattice-only decay without continuum spectral meaning.
\end{itemize}

\paragraph{What was learned physically}

\begin{itemize}
\item The analysis should now advance by testing typed bridge lemmas, not by accumulating more route families.
\item A strategy only counts as learning if it changes the mathematical dependency order: Hamiltonian claims move behind reconstruction gates, decay claims move behind uniformity/density gates, and RG claims move behind OS/nontriviality gates.
\item The next useful artifact is therefore a bridge-lemma packet with definitions, assumptions, equations, failure tests, and conditional skill triggers.
\end{itemize}

\paragraph{Decision}

The next iteration should produce explicit bridge-lemma packets, not another broad slogan.

\subsection*{We then added reductio branches to test whether the negated worlds create real mathematical pressure.}

\paragraph{Physics mechanism}

A reductio is physically useful only when the negated world leaves a detectable trace. If there is no gap, there should be a sequence of vacuum-orthogonal gauge-invariant states with energies tending to zero, or equivalently spectral weight accumulating at the origin for some determining observable family. If OS positivity fails, the Euclidean theory cannot reconstruct a positive-norm quantum Hilbert space. If the continuum limit is trivial, connected gauge-invariant cumulants collapse. These are physical failure modes, not just logical negations.

\paragraph{What did not pay off or remains blocked}

\begin{itemize}
\item No contradiction branch currently proves Yang-Mills; each branch is only promotable as a proof obligation.
\item The no-gap reductio still depends on observable density and cutoff-uniform decay.
\item The OS and triviality reductios still depend on closedness of reflection positivity and a same-theory nontriviality witness.
\end{itemize}

\paragraph{What was learned physically}

\begin{itemize}
\item The loop learned to use contradiction as a disciplined pressure test rather than an excuse to skip construction.
\item A branch succeeds only if the contradiction target is accepted or validator-backed in the same theory.
\item A failed branch is still valuable when it tells us whether to pivot toward density, closedness, nontriviality, or uniform-decay evidence.
\end{itemize}

\paragraph{Decision}

Keep all three reductio branches active as proof-obligation pressure tests, not as proof results.

\section*{Gedanken Experiment Release Summary}

\begin{itemize}
\item Packet count: 16
\item Live packet count: 16
\item Queue handoff count: 32
\item Quarantine count: 0
\item Invalid packet count: 0
\item Proof-safety status: conceptual\_probes\_only
\item Gedanken packets may affect branch control only through validated follow-up queue items, transfer verdicts, and qualitative graph annotations. They are not counted as proof progress unless a later derivation, computation, or source-backed certificate discharges the relevant obligation.
\end{itemize}

\end{document}
